\section{位相の定義}

\subsection{開集合}

\begin{framed}
    \begin{dfn}[位相の定義と開集合]\label{dfn::topology}
        $S$を空でない集合とする.$2^S$の部分集合$\mathcal{O}$が以下の条件を満たすとき,$\mathcal{O}$を$S$の\textbf{位相}であるといい,位相$\mathcal{O}$が与えられた集合$S$を\textbf{位相空間}という.

        \begin{enumerate}
            \item $\displaystyle\emptyset,S\in\mathcal{O}.$
            \item $O_1,O_2\in\mathcal{O}$ならば$O_1\cap O_2\in\mathcal{O}.$
            \item 任意の添字$\lambda\in\Lambda$において$O_\lambda\in\mathcal{O}$ならば$\displaystyle\bigcup_{\lambda\in\Lambda}O_\lambda\in\mathcal{O}.$
        \end{enumerate}

        $S$を$\textbf{台}$(または\textbf{台集合})といい,$S$の元を\textbf{点}という.
        また,$\mathcal{O}$に属する$S$の部分集合を\textbf{$\mathcal{O}$-開集合}(あるいは単に\textbf{開集合})という.

    \end{dfn}
\end{framed}

位相空間は$(S,\mathcal{O})$のように書くが,文脈から与えられている位相が明らかな場合は単に$S$と書くこともある.

\begin{exa}
    $\mathcal{O}=\{\emptyset,S\}$は$S$の位相の条件を満たす.この位相を\textbf{密着位相}といい,位相空間$(S,\{\emptyset,S\})$を\textbf{密着空間}という.
\end{exa}

\begin{exa}
    $\mathcal{O}=2^S$は$S$の位相の条件を満たす.この位相を\textbf{離散位相}といい,位相空間$(S,2^S)$を\textbf{離散空間}という.
\end{exa}

\begin{exa}
    $(S,\mathcal{O})$を位相空間,$A$を空でない集合とする.
    $\mathcal{O}_A=\{A\cap O\mid O\in\mathcal{O}\}$は$S$の位相の条件を満たす.この位相を集合$A$上の$\mathcal{O}$に関する\textbf{相対位相}といい,位相空間$(A,\mathcal{O}_A)$を$(S,\mathcal{O})$の\textbf{部分空間}という.

\end{exa}


\subsection{閉集合}

\begin{framed}
    \begin{dfn}[位相の定義と開集合]\label{dfn::closed-set}
        $\mathcal{O}$-開集合の補集合を\textbf{$\mathcal{O}$-閉集合}(あるいは単に\textbf{閉集合})という.
    \end{dfn}
\end{framed}

言い換えるならば,$\mathcal{O}$-閉集合全体の集合は$\{O^c\mid O\in\mathcal{O}\}$で表せられる.

\begin{framed}
    \begin{thm}[閉集合]\label{thm::closed-set}
        $(S,\mathcal{O})$を位相空間とする.$S$の閉集合全体の集合を$\mathcal{A}$としたとき,$\mathcal{A}$は次を満たす.

        \begin{enumerate}
            \item$\displaystyle\emptyset,S\in\mathcal{A}.$
            \item$A_1,A_2\in\mathcal{A}$ならば$A_1\cup A_2\in\mathcal{A}.$
            \item 任意の添字$\lambda\in\Lambda$において$A_\lambda\in\mathcal{A}$ならば$\displaystyle\bigcap_{\lambda\in\Lambda}A_\lambda\in\mathcal{A}.$
        \end{enumerate}
    \end{thm}
\end{framed}

\begin{proof}
    それぞれ開集合の定義とde Morganの法則から直ちに導かれる.例えば(3)は次のように示される.

    $A_\lambda^c=O_\lambda$とすれば,定義\ref{dfn::topology}(3)より
    $$\displaystyle\bigcup_{\lambda\in\Lambda}{A^c}=\left(\bigcap_{\lambda\in\Lambda}{A}\right)^c\in\mathcal{O}$$
    である.よって閉集合の定義より$\displaystyle\bigcap_{\lambda\in\Lambda}A_\lambda\in\mathcal{A}.$
\end{proof}

逆に,定理\ref{thm::closed-set}の3条件を満たす集合族$\mathcal{A}$が与えられたならば,$\mathcal{A}$が$(S,\mathcal{O})$の閉集合族であるような位相$\mathcal{O}$は一意に定まる.
すなわち,$S$を台とする位相空間は開集合族$\mathcal{O}$の代わりに閉集合族$\mathcal{A}$を定めても決定できる.

\begin{rem*}
    開集合と閉集合は補集合の意味で双対の関係にあるが,開集合でない集合が閉集合,あるいは閉集合でない集合が開集合であるとは限らない.

    例えば,任意の位相空間$(S,\mathcal{O})$において空集合$\emptyset$とその台集合$S$は開かつ閉集合である.
    また,離散空間のとき$\emptyset$と$S$を除く$S$の任意の部分集合は閉集合でも開集合でもない.
\end{rem*}

\subsection{内部・開核作用素}

\begin{framed}
    \begin{dfn}[]\label{dfn::inner}
        $(S,\mathcal{O})$ を位相空間とする. $S$ の部分集合 $A$ において, $$\bigcup\{O\in\mathcal{O} \mid O\subset A\}$$ となる集合を $A$ の \textbf{内部} あるいは \textbf{開核} といい, $A^\circ$ で表す.
    \end{dfn}
\end{framed}


\begin{framed}
    \begin{dfn}[]\label{dfn::inner-op}
        $(S,\mathcal{O})$ を位相空間, $A,B$ を $S$ の任意の部分集合とする. 次の条件を満たす写像 $i$ を \textbf{閉核作用素} という.

        \begin{enumerate}
            \item $i(S)=S.$
            \item $i(A)\subset A.$
            \item $i(A\cap B)=i(A)\cap i(B).$
            \item $(i\circ i)(A)=i(A).$
        \end{enumerate}
    \end{dfn}
\end{framed}

\begin{framed}
    \begin{thm}[]\label{thm::inner-op}
        定義\ref{dfn::inner}で定めた内部は定義\ref{dfn::inner-op}での4条件を満たす.
    \end{thm}
\end{framed}



\begin{framed}
    \begin{dfn}[]\label{dfn::closure-op}
        $(S,\mathcal{O})$ を位相空間, $A,B$ を $S$ の任意の部分集合とする. 次の条件を満たす演算 $\overline{\cdot}$ を \textbf{閉包作用素} という.

        \begin{enumerate}
            \item $\overline{\emptyset}=\emptyset.$
            \item $\overline{A}\supset A.$
            \item $\overline{A\cup B}=\overline{A}\cup\overline{B}.$
            \item $\overline{\overline{A}}=\overline{A}.$
        \end{enumerate}
    \end{dfn}
\end{framed}